\begin{proposition}[The structure-sheaf input in the precise form used below]\label{inv:A1.0}\label{wb:thm:structure-direct-images-audit}\label{wb:lem:cusp-structure-sheaf-audit}
This is the structure-sheaf calculation of Claude and Alp\"oge: \cite[Lemma~9.12, pp.~72--74; Proposition~9.13, pp.~74--76, and the base-change clarification on p.~92]{AC2026}. The direct images $R^qf_*\OO_X$, for $q=0,1,2$, are
$\OO$, $G=\OO\gamma\oplus\OO(-1)$, and $\OO(-1)$.
The inclusion $\OO\gamma\subset G$ is a subbundle.  At the cusp,
$(\gamma,u)$ is a holomorphic frame, and cup product identifies
$\bigwedge^2G$ with $R^2f_*\OO_X$.  In the cyclic covering coordinate
$s$, the action on $G$ has characters $(1,\zeta^{-a})$, where
$(m,a)=(3,2)$ or $(4,1)$ and $\zeta=e^{-2\pi i/m}$.
These assertions concern full scheme fibres, including $3S_1$ and $4S_2$.
At the cusp, the natural map $H^1(W;\mathbb C)\to H^1(W,\mathcal O_W)$ is an isomorphism. In degree two the comparison is surjective and is injective on $\ker[H^2(W;\mathbb C)\to H^2(\mathrm{dP}_6;\mathbb C)]$. The cup product $\bigwedge^2H^1(W,\mathcal O_W)\to H^2(W,\mathcal O_W)$ is an isomorphism. Every base-change map $(R^qf_*\mathcal O_X)\otimes k(b)\to H^q(X_b,\mathcal O_{X_b})$ is an isomorphism, including at the nonreduced fibres.
\end{proposition}
\Uses{\cite[Lemma~9.12, pp.~72--74; Proposition~9.13, pp.~74--76; p.~92]{AC2026}}
